% Source: book pp. 243--250 (PDF 257--264). No solutions in the manual.
% a numbered display (Axler numbers a few equations with the item counter):
%   text\axeqstep{ax:7.28}  \[ ... \tag{\thetheorem} \]
\DeclareRobustCommand\axeqstep[1]{\refstepcounter{theorem}\label{#1}}
% the small italic comment that follows some exercises
\section{Spectral Theorem}

Recall that a diagonal matrix is a square matrix that is $0$ everywhere except
possibly on the diagonal. Recall that an operator on $V$ is called
diagonalizable if the operator has a diagonal matrix with respect to some basis
of $V$. Recall also that this happens if and only if there is a basis of $V$
consisting of eigenvectors of the operator (see \ref{ax:5.55}).

The nicest operators on $V$ are those for which there is an \emph{orthonormal}
basis of $V$ with respect to which the operator has a diagonal matrix. These are
precisely the operators $T\in\Lin(V)$ such that there is an orthonormal basis of
$V$ consisting of eigenvectors of $T$. Our goal in this section is to prove the
spectral theorem, which characterizes these operators as the self-adjoint
operators when $\F=\R$ and as the normal operators when $\F=\C$.

The spectral theorem is probably the most useful tool in the study of operators
on inner product spaces. Its extension to certain infinite-dimensional inner
product spaces (see, for example, Section~10D of the author's book
\emph{Measure, Integration \& Real Analysis}) plays a key role in functional
analysis.

Because the conclusion of the spectral theorem depends on $\F$, we will break
the spectral theorem into two pieces, called the real spectral theorem and the
complex spectral theorem.

\subsection{Real Spectral Theorem}

To prove the real spectral theorem, we will need two preliminary results. These
preliminary results hold on both real and complex inner product spaces, but
they are not needed for the proof of the complex spectral theorem.

\marginnote{This completing-the-square technique can be used to derive the
quadratic formula.}%
You could guess that the next result is true and even discover its proof by
thinking about quadratic polynomials with real coefficients. Specifically,
suppose $b,c\in\R$ and $b^2<4c$. Let $x$ be a real number. Then
\[ x^2+bx+c = \Bigl(x+\frac{b}{2}\Bigr)^2+\Bigl(c-\frac{b^2}{4}\Bigr) > 0. \]
In particular, $x^2+bx+c$ is an invertible real number (a convoluted way of
saying that it is not $0$). Replacing the real number $x$ with a self-adjoint
operator (recall the analogy between real numbers and self-adjoint operators)
leads to the next result.

\begin{result}[Invertible quadratic expressions]\label{ax:7.26}
Suppose $T\in\Lin(V)$ is self-adjoint and $b,c\in\R$ are such that $b^2<4c$.
Then
\[ T^2+bT+cI \]
is an invertible operator.
\end{result}

\begin{proof}
Let $v$ be a nonzero vector in $V$. Then
\begin{align*}
\bigl\langle (T^2+bT+cI)v,v\bigr\rangle
  &= \ip{T^2v}{v}+b\ip{Tv}{v}+c\ip{v}{v}\\
  &= \ip{Tv}{Tv}+b\ip{Tv}{v}+c\norm{v}^2\\
  &\ge \norm{Tv}^2-\abs{b}\,\norm{Tv}\,\norm{v}+c\norm{v}^2\\
  &= \Bigl(\norm{Tv}-\frac{\abs{b}\,\norm{v}}{2}\Bigr)^2
     +\Bigl(c-\frac{b^2}{4}\Bigr)\norm{v}^2\\
  &> 0,
\end{align*}
where the third line above holds by the Cauchy--Schwarz inequality
(\ref{ax:6.14}). The last inequality implies that $(T^2+bT+cI)v\ne0$. Thus
$T^2+bT+cI$ is injective, which implies that it is invertible (see
\ref{ax:3.65}).
\end{proof}

The next result will be a key tool in our proof of the real spectral theorem.

\begin{result}[Minimal polynomial of self-adjoint operator]\label{ax:7.27}
Suppose $T\in\Lin(V)$ is self-adjoint. Then the minimal polynomial of $T$
equals $(z-\lambda_1)\cdots(z-\lambda_m)$ for some
$\lambda_1,\dots,\lambda_m\in\R$.
\end{result}

\begin{proof}
First suppose $\F=\C$. The zeros of the minimal polynomial of $T$ are the
eigenvalues of $T$ [by \ref{ax:5.27}(a)]. All eigenvalues of $T$ are real (by
\ref{ax:7.12}). Thus the second version of the fundamental theorem of algebra
(see \ref{ax:4.13}) tells us that the minimal polynomial of $T$ has the desired
form.

Now suppose $\F=\R$. By the factorization of a polynomial over $\R$ (see
\ref{ax:4.16}) there exist $\lambda_1,\dots,\lambda_m\in\R$ and
$b_1,\dots,b_N,c_1,\dots,c_N\in\R$ with $b_k^2<4c_k$ for each $k$ such that the
minimal polynomial of $T$ equals\axeqstep{ax:7.28}
\[ (z-\lambda_1)\cdots(z-\lambda_m)(z^2+b_1z+c_1)\cdots(z^2+b_Nz+c_N);
   \tag{\thetheorem} \]
here either $m$ or $N$ might equal $0$, meaning that there are no terms of the
corresponding form. Now
\[ (T-\lambda_1I)\cdots(T-\lambda_mI)(T^2+b_1T+c_1I)\cdots(T^2+b_NT+c_NI) = 0. \]
If $N>0$, then we could multiply both sides of the equation above on the right
by the inverse of $T^2+b_NT+c_NI$ (which is an invertible operator by
\ref{ax:7.26}) to obtain a polynomial expression of $T$ that equals $0$. The
corresponding polynomial would have degree two less than the degree of
\ref{ax:7.28}, violating the minimality of the degree of the polynomial with
this property. Thus we must have $N=0$, which means that the minimal polynomial
in \ref{ax:7.28} has the form $(z-\lambda_1)\cdots(z-\lambda_m)$, as desired.
\end{proof}

The result above along with \ref{ax:5.27}(a) implies that every self-adjoint
operator has an eigenvalue. In fact, as we will see in the next result,
self-adjoint operators have enough eigenvectors to form a basis.

The next result, which gives a complete description of the self-adjoint
operators on a real inner product space, is one of the major theorems in linear
algebra.

\begin{result}[Real spectral theorem]\label{ax:7.29}
Suppose $\F=\R$ and $T\in\Lin(V)$. Then the following are equivalent.
\begin{parts}
\item $T$ is self-adjoint.
\item $T$ has a diagonal matrix with respect to some orthonormal basis of $V$.
\item $V$ has an orthonormal basis consisting of eigenvectors of $T$.
\end{parts}
\end{result}

\begin{proof}
First suppose (a) holds, so $T$ is self-adjoint. Our results on minimal
polynomials, specifically \ref{ax:6.37} and \ref{ax:7.27}, imply that $T$ has
an upper-triangular matrix with respect to some orthonormal basis of $V$. With
respect to this orthonormal basis, the matrix of $T^*$ is the transpose of the
matrix of $T$. However, $T^*=T$. Thus the transpose of the matrix of $T$ equals
the matrix of $T$. Because the matrix of $T$ is upper-triangular, this means
that all entries of the matrix above and below the diagonal are $0$. Hence the
matrix of $T$ is a diagonal matrix with respect to the orthonormal basis. Thus
(a) implies (b).

Conversely, now suppose (b) holds, so $T$ has a diagonal matrix with respect to
some orthonormal basis of $V$. That diagonal matrix equals its transpose. Thus
with respect to that basis, the matrix of $T^*$ equals the matrix of $T$. Hence
$T^*=T$, proving that (b) implies (a).

The equivalence of (b) and (c) follows from the definitions [or see the proof
that (a) and (b) are equivalent in \ref{ax:5.55}].
\end{proof}

\begin{example}[An orthonormal basis of eigenvectors for an operator]%
\label{ax:7.30}
Consider the operator $T$ on $\R^3$ whose matrix (with respect to the standard
basis) is
\[ \begin{pmatrix} 14 & -13 & 8\\ -13 & 14 & 8\\ 8 & 8 & -7 \end{pmatrix}. \]
This matrix with real entries equals its transpose; thus $T$ is self-adjoint.
As you can verify,
\[ \frac{(1,-1,0)}{\sqrt2},\ \frac{(1,1,1)}{\sqrt3},\ \frac{(1,1,-2)}{\sqrt6} \]
is an orthonormal basis of $\R^3$ consisting of eigenvectors of $T$. With
respect to this basis, the matrix of $T$ is the diagonal matrix
\[ \begin{pmatrix} 27 & 0 & 0\\ 0 & 9 & 0\\ 0 & 0 & -15 \end{pmatrix}. \]
\end{example}

See Exercise~17 for a version of the real spectral theorem that applies
simultaneously to more than one operator.

\subsection{Complex Spectral Theorem}

The next result gives a complete description of the normal operators on a
complex inner product space.

\begin{result}[Complex spectral theorem]\label{ax:7.31}
Suppose $\F=\C$ and $T\in\Lin(V)$. Then the following are equivalent.
\begin{parts}
\item $T$ is normal.
\item $T$ has a diagonal matrix with respect to some orthonormal basis of $V$.
\item $V$ has an orthonormal basis consisting of eigenvectors of $T$.
\end{parts}
\end{result}

\begin{proof}
First suppose (a) holds, so $T$ is normal. By Schur's theorem (\ref{ax:6.38}),
there is an orthonormal basis $e_1,\dots,e_n$ of $V$ with respect to which $T$
has an upper-triangular matrix. Thus we can write\axeqstep{ax:7.32}
\[ \Mat\bigl(T,(e_1,\dots,e_n)\bigr) =
   \begin{pmatrix} a_{1,1} & \cdots & a_{1,n}\\ & \ddots & \vdots\\
   0 & & a_{n,n} \end{pmatrix}. \tag{\thetheorem} \]
We will show that this matrix is actually a diagonal matrix.

We see from the matrix above that
\begin{align*}
\norm{Te_1}^2 &= \abs{a_{1,1}}^2,\\
\norm{T^*e_1}^2 &= \abs{a_{1,1}}^2+\abs{a_{1,2}}^2+\cdots+\abs{a_{1,n}}^2.
\end{align*}

Because $T$ is normal, $\norm{Te_1}=\norm{T^*e_1}$ (see \ref{ax:7.20}). Thus the
two equations above imply that all entries in the first row of the matrix in
\ref{ax:7.32}, except possibly the first entry $a_{1,1}$, equal $0$.

Now \ref{ax:7.32} implies
\[ \norm{Te_2}^2 = \abs{a_{2,2}}^2 \]
(because $a_{1,2}=0$, as we showed in the paragraph above) and
\[ \norm{T^*e_2}^2 = \abs{a_{2,2}}^2+\abs{a_{2,3}}^2+\cdots+\abs{a_{2,n}}^2. \]
Because $T$ is normal, $\norm{Te_2}=\norm{T^*e_2}$. Thus the two equations
above imply that all entries in the second row of the matrix in \ref{ax:7.32},
except possibly the diagonal entry $a_{2,2}$, equal $0$.

Continuing in this fashion, we see that all nondiagonal entries in the matrix
\ref{ax:7.32} equal $0$. Thus (b) holds, completing the proof that (a) implies
(b).

Now suppose (b) holds, so $T$ has a diagonal matrix with respect to some
orthonormal basis of $V$. The matrix of $T^*$ (with respect to the same basis)
is obtained by taking the conjugate transpose of the matrix of $T$; hence $T^*$
also has a diagonal matrix. Any two diagonal matrices commute; thus $T$
commutes with $T^*$, which means that $T$ is normal. In other words, (a) holds,
completing the proof that (b) implies (a).

The equivalence of (b) and (c) follows from the definitions (also see
\ref{ax:5.55}).
\end{proof}

See Exercises~13 and~20 for alternative proofs that (a) implies (b) in the
previous result.

Exercises~14 and~15 interpret the real spectral theorem and the complex
spectral theorem by expressing the domain space as an orthogonal direct sum of
eigenspaces.

See Exercise~16 for a version of the complex spectral theorem that applies
simultaneously to more than one operator.

The main conclusion of the complex spectral theorem is that every normal
operator on a complex finite-dimensional inner product space is diagonalizable
by an orthonormal basis, as illustrated by the next example.

\begin{example}[An orthonormal basis of eigenvectors for an operator]%
\label{ax:7.33}
Consider the operator $T\in\Lin(\C^2)$ defined by $T(w,z)=(2w-3z,3w+2z)$. The
matrix of $T$ (with respect to the standard basis) is
\[ \begin{pmatrix} 2 & -3\\ 3 & 2 \end{pmatrix}. \]
As we saw in Example~\ref{ax:7.19}, $T$ is a normal operator.

As you can verify,
\[ (1/\sqrt2)(i,1),\ (1/\sqrt2)(-i,1) \]
is an orthonormal basis of $\C^2$ consisting of eigenvectors of $T$, and with
respect to this basis the matrix of $T$ is the diagonal matrix
\[ \begin{pmatrix} 2+3i & 0\\ 0 & 2-3i \end{pmatrix}. \]
\end{example}

% ------------------------------------------------------------------ exercises
\begin{exc}

\begin{exercise}
Prove that a normal operator on a complex inner product space is self-adjoint
if and only if all its eigenvalues are real.
\exercisenote{This exercise strengthens the analogy (for normal operators)
between self-adjoint operators and real numbers.}
\end{exercise}

\begin{exercise}
Suppose $\F=\C$. Suppose $T\in\Lin(V)$ is normal and has only one eigenvalue.
Prove that $T$ is a scalar multiple of the identity operator.
\end{exercise}

\begin{exercise}
Suppose $\F=\C$ and $T\in\Lin(V)$ is normal. Prove that the set of eigenvalues
of $T$ is contained in $\{0,1\}$ if and only if there is a subspace $U$ of $V$
such that $T=P_U$.
\end{exercise}

\begin{exercise}
Prove that a normal operator on a complex inner product space is skew (meaning
it equals the negative of its adjoint) if and only if all its eigenvalues are
purely imaginary (meaning that they have real part equal to $0$).
\end{exercise}

\begin{exercise}
Prove or give a counterexample: If $T\in\Lin(\C^3)$ is a diagonalizable
operator, then $T$ is normal (with respect to the usual inner product).
\end{exercise}

\begin{exercise}
Suppose $V$ is a complex inner product space and $T\in\Lin(V)$ is a normal
operator such that $T^9=T^8$. Prove that $T$ is self-adjoint and $T^2=T$.
\end{exercise}

\begin{exercise}
Give an example of an operator $T$ on a complex vector space such that
$T^9=T^8$ but $T^2\ne T$.
\end{exercise}

\begin{exercise}
Suppose $\F=\C$ and $T\in\Lin(V)$. Prove that $T$ is normal if and only if
every eigenvector of $T$ is also an eigenvector of $T^*$.
\end{exercise}

\begin{exercise}
Suppose $\F=\C$ and $T\in\Lin(V)$. Prove that $T$ is normal if and only if
there exists a polynomial $p\in\Poly(\C)$ such that $T^*=p(T)$.
\end{exercise}

\begin{exercise}
Suppose $V$ is a complex inner product space. Prove that every normal operator
on $V$ has a square root.
\exercisenote{An operator $S\in\Lin(V)$ is called a \textbf{square root} of
$T\in\Lin(V)$ if $S^2=T$. We will discuss more about square roots of operators
in Sections~7C and~8C.}
\end{exercise}

\begin{exercise}
Prove that every self-adjoint operator on $V$ has a cube root.
\exercisenote{An operator $S\in\Lin(V)$ is called a \textbf{cube root} of
$T\in\Lin(V)$ if $S^3=T$.}
\end{exercise}

\begin{exercise}
Suppose $V$ is a complex vector space and $T\in\Lin(V)$ is normal. Prove that
if $S$ is an operator on $V$ that commutes with $T$, then $S$ commutes with
$T^*$.
\exercisenote{The result in this exercise is called Fuglede's theorem.}
\end{exercise}

\begin{exercise}
Without using the complex spectral theorem, use the version of Schur's theorem
that applies to two commuting operators (take $\mathcal{E}=\{T,T^*\}$ in
Exercise~20 in Section~6B) to give a different proof that if $\F=\C$ and
$T\in\Lin(V)$ is normal, then $T$ has a diagonal matrix with respect to some
orthonormal basis of $V$.
\end{exercise}

\begin{exercise}
Suppose $\F=\R$ and $T\in\Lin(V)$. Prove that $T$ is self-adjoint if and only
if all pairs of eigenvectors corresponding to distinct eigenvalues of $T$ are
orthogonal and $V=E(\lambda_1,T)\oplus\cdots\oplus E(\lambda_m,T)$, where
$\lambda_1,\dots,\lambda_m$ denote the distinct eigenvalues of $T$.
\end{exercise}

\begin{exercise}
Suppose $\F=\C$ and $T\in\Lin(V)$. Prove that $T$ is normal if and only if all
pairs of eigenvectors corresponding to distinct eigenvalues of $T$ are
orthogonal and $V=E(\lambda_1,T)\oplus\cdots\oplus E(\lambda_m,T)$, where
$\lambda_1,\dots,\lambda_m$ denote the distinct eigenvalues of $T$.
\end{exercise}

\begin{exercise}
Suppose $\F=\C$ and $\mathcal{E}\subseteq\Lin(V)$. Prove that there is an
orthonormal basis of $V$ with respect to which every element of $\mathcal{E}$
has a diagonal matrix if and only if $S$ and $T$ are commuting normal operators
for all $S,T\in\mathcal{E}$.
\exercisenote{This exercise extends the complex spectral theorem to the context
of a collection of commuting normal operators.}
\end{exercise}

\begin{exercise}
Suppose $\F=\R$ and $\mathcal{E}\subseteq\Lin(V)$. Prove that there is an
orthonormal basis of $V$ with respect to which every element of $\mathcal{E}$
has a diagonal matrix if and only if $S$ and $T$ are commuting self-adjoint
operators for all $S,T\in\mathcal{E}$.
\exercisenote{This exercise extends the real spectral theorem to the context of
a collection of commuting self-adjoint operators.}
\end{exercise}

\begin{exercise}
Give an example of a real inner product space $V$, an operator $T\in\Lin(V)$,
and real numbers $b,c$ with $b^2<4c$ such that
\[ T^2+bT+cI \]
is not invertible.
\exercisenote{This exercise shows that the hypothesis that $T$ is self-adjoint
cannot be deleted in \ref{ax:7.26}, even for real vector spaces.}
\end{exercise}

\begin{exercise}
Suppose $T\in\Lin(V)$ is self-adjoint and $U$ is a subspace of $V$ that is
invariant under $T$.
\begin{parts}
\item Prove that $U^\perp$ is invariant under $T$.
\item Prove that $T|_U\in\Lin(U)$ is self-adjoint.
\item Prove that $T|_{U^\perp}\in\Lin(U^\perp)$ is self-adjoint.
\end{parts}
\end{exercise}

\begin{exercise}
Suppose $T\in\Lin(V)$ is normal and $U$ is a subspace of $V$ that is invariant
under $T$.
\begin{parts}
\item Prove that $U^\perp$ is invariant under $T$.
\item Prove that $U$ is invariant under $T^*$.
\item Prove that $(T|_U)^*=(T^*)|_U$.
\item Prove that $T|_U\in\Lin(U)$ and $T|_{U^\perp}\in\Lin(U^\perp)$ are normal
  operators.
\end{parts}
\exercisenote{This exercise can be used to give yet another proof of the
complex spectral theorem (use induction on $\dim V$ and the result that $T$ has
an eigenvector).}
\end{exercise}

\begin{exercise}
Suppose that $T$ is a self-adjoint operator on a finite-dimensional inner
product space and that $2$ and $3$ are the only eigenvalues of $T$. Prove that
\[ T^2-5T+6I = 0. \]
\end{exercise}

\begin{exercise}
Give an example of an operator $T\in\Lin(\C^3)$ such that $2$ and $3$ are the
only eigenvalues of $T$ and $T^2-5T+6I\ne0$.
\end{exercise}

\begin{exercise}
Suppose $T\in\Lin(V)$ is self-adjoint, $\lambda\in\F$, and $\epsilon>0$.
Suppose there exists $v\in V$ such that $\norm{v}=1$ and
\[ \norm{Tv-\lambda v} < \epsilon. \]
Prove that $T$ has an eigenvalue $\lambda'$ such that
$\abs{\lambda-\lambda'}<\epsilon$.
\exercisenote{This exercise shows that for a self-adjoint operator, a number
that is close to satisfying an equation that would make it an eigenvalue is
close to an eigenvalue.}
\end{exercise}

\begin{exercise}
Suppose $U$ is a finite-dimensional vector space and $T\in\Lin(U)$.
\begin{parts}
\item Suppose $\F=\R$. Prove that $T$ is diagonalizable if and only if there is
  a basis of $U$ such that the matrix of $T$ with respect to this basis equals
  its transpose.
\item Suppose $\F=\C$. Prove that $T$ is diagonalizable if and only if there is
  a basis of $U$ such that the matrix of $T$ with respect to this basis
  commutes with its conjugate transpose.
\end{parts}
\exercisenote{This exercise adds another equivalence to the list of conditions
equivalent to diagonalizability in \ref{ax:5.55}.}
\end{exercise}

\begin{exercise}
Suppose that $T\in\Lin(V)$ and there is an orthonormal basis $e_1,\dots,e_n$ of
$V$ consisting of eigenvectors of $T$, with corresponding eigenvalues
$\lambda_1,\dots,\lambda_n$. Show that if $k\in\{1,\dots,n\}$, then the
pseudoinverse $T^\dagger$ satisfies the equation
\[ T^\dagger e_k = \begin{cases} \frac{1}{\lambda_k}e_k & \text{if }
   \lambda_k\ne0,\\ 0 & \text{if } \lambda_k=0. \end{cases} \]
\end{exercise}
\end{exc}
